% ============================================================
% kbs_data.tex — KBS journal Data & corpus section (KE voice).
% Section label: sec:data (referenced by the Introduction draft).
% Conference-scale numbers shown; journal extension to a 1,300-book
% corpus marked \TODO{E1}; panel-validated gold marked \TODO{E2}.
% Bibliography: kbs_ke_references.bib + paper/references.bib.
% ============================================================

\section{Data and Ground Truth}
\label{sec:data}

The knowledge-engineering claims in this paper require a collection in
which (i) subject knowledge is professionally available for a subset of
works, (ii) the full text is available so that acquisition can be studied
with and without content conditioning, and (iii) a second professional
source exists to quantify inter-source disagreement in the gold itself. We
build the collection from three open sources, all public-domain or openly
licensed: Project Gutenberg full text (with LC-derived subject metadata),
OpenLibrary MARC-derived metadata (title, author, Dewey classification),
and HathiTrust bibliographic records as a professional cross-check of the
gold, alongside other open bibliographic
corpora~\cite{hathitrust_ef, blbooks}. This section describes the corpus, the gold knowledge base, the
validation of the gold, and the released benchmark instruments.

\subsection{Corpus}
\label{sec:data:corpus}

The retrieval and acquisition experiments run over a curated collection of
English-language public-domain books. Works are joined across Gutenberg and
OpenLibrary by fuzzy title/author matching: of 2{,}000 sampled books,
1{,}359 (68\%) matched to a work record carrying subject knowledge
\TODO{E1: scale the corpus from 227 to ~1,300 books; report the final join
and chunk counts}. For each joined work the collection stores the full
text, chunked for retrieval (conference scale: 227 books, 88{,}264
chunks), and the catalog record---identity fields (title, author, year) and
subject knowledge (headings, classification)---which is the unit the
acquisition component operates on and the knowledge-based system retrieves.

Two properties of the corpus matter for the knowledge claims. First, the
works are \emph{well-known public-domain literature}, which bounds the
memorization concern: a large model may recall the real headings of famous
works rather than perform knowledge acquisition, so we report a
popularity-tier analysis that compares acquisition quality across download
counts to bound recall effects. Second, subject knowledge is
\emph{unevenly distributed}---not every record carries headings or
classification---which is exactly the knowledge-base-completeness
heterogeneity that the paper exploits. Records that carry no subject
knowledge are the population the acquisition component is asked to fill.

\subsection{Knowledge base and gold}
\label{sec:data:gold}

The gold knowledge base covers \TODO{E1: extend} records with subject
knowledge in both target representations:

\textbf{Subject headings (thesaurus).} Ground-truth headings are
LC-derived Gutenberg subject metadata, i.e., headings in Library of Congress
Subject Headings form. Every record in the benchmark carries at least one
gold heading (conference scale: 600 records, 1{,}405 gold headings).

\textbf{Classification (taxonomy).} Ground-truth Dewey numbers come from
OpenLibrary's aggregated MARC 082 field. Because OpenLibrary aggregates
editions, a record may carry several valid Dewey numbers; classification
gold is therefore scored as \emph{any-gold-match} (conference scale: 204 of
600 records carry gold Dewey).

\subsection{Validating the gold itself}
\label{sec:data:gold-validation}

The gold must be trustworthy for validation claims, so we cross-check it
against an independent professional source---HathiTrust MARC---on a
subsample (conference scale: 102 of 224 pilot books found in HathiTrust; 36
carried LCSH fields and 19 carried Dewey). The cross-check is
methodologically decisive for how acquisition quality must be scored:
where both sources carry Dewey, 17/19 books agree at the 3-digit class
(89\%), while exact subject-heading agreement between the two professional
sources is far lower (9 of 34 books share an exact heading; mean recall of
one source's headings in the other is 15\%). Two professional cataloging
sources therefore disagree on exact heading form most of the time. The
multi-level rubric (exact / semantic-equivalence / acceptable) exists
precisely so that inter-cataloger variation is not scored as acquisition
error; and the \emph{agreement ceiling}---the rate at which two
professionals agree---is the correct reference for judging LLM acquisition,
not exact match \TODO{E2: measure the inter-cataloger agreement ceiling on
a panel-validated sample and report LLM scores against it}. The HathiTrust
cross-check also provides the second professional source that motivates the
decomposition of valid-but-unmatched predicted headings
(Section~\ref{sec:valid}).

\subsection{Benchmark instruments}
\label{sec:data:benchmark}

Two released instrument sets operationalize the two sides of the
knowledge-engineering cycle.

\textbf{Discovery questions (utilization side).} A released set of
discovery questions \TODO{E1: extend} in three classes---known-item
(locate a work by title/author), topical (find a work about a subject), and
bibliographic-fact (report a record field such as headings or class)---each
with gold answers linked to record identifiers so correctness is verifiable
by construction. Question wording is LLM-polished into natural patron
language while gold answers remain frozen. Of the released set,
\TODO{E1: extend} fall inside the evaluation corpus and form the retrieval
pool; the topical subset (conference scale: 338 of 1{,}076 in-corpus
questions) drives the headline claims and is the subject of the
single-gold-bias analysis in Section~\ref{sec:util}.

\textbf{Cataloging records (validation side).} The gold knowledge base
above, released with the multi-level scoring protocol, the error taxonomy,
and live authority checking against \texttt{id.loc.gov}. Authority
validation is applied as a \emph{measurement} on predicted headings---it
classifies them into the taxonomy (authority-violating vs.\ valid-but-
unmatched)---rather than as a gate that filters what is indexed; we do not
test a gated variant and claim nothing for one.

\subsection{Data provenance and reuse}
\label{sec:data:provenance}

All sources are open and license-compatible: Project Gutenberg (public
domain), OpenLibrary (CC0/MARC-derived), HathiTrust public bibliographic
API, and the Library of Congress authorities via \texttt{id.loc.gov}.
Corpus construction, joining, and indexing are fully scripted and released;
every number reported in this paper traces to a committed JSON artifact,
and the benchmark instruments are released for independent reuse
\TODO{E1: update released corpus and question counts to the journal scale}.
No patron or circulation data is used anywhere; institutional review was
confirmed not to be required for the public-domain collection, and the
cataloger-panel study follows the protocol in
Section~\ref{sec:valid} \TODO{E2: attach panel consent/IRB note}.
